%============================================================ 초록
\begin{abstract}
\noindent
저비용 자폭 무인수상정(USV) 군집은 재래식 수단으로 대응하기 어렵다. 전자전에 의한 연성
대응은 유선 조종 표적에 닿지 않고, 경성 요격은 표적보다 비싼 유도탄을 한 번에 하나씩
소모한다. 본 연구는 대안으로 물리적 포획을 채택하여, 적보다 느린 소수의 방어정이 모선의
접근 회랑을 그물벽으로 미리 차단하는 방어 체계를 제안한다. 체계는 같은 주기로 갱신되지만
\emph{응답 지연이 크게 다른} 두 계층으로 구성되며, 느린 계층을 비동기로 분리하여 그 지연이
제어 루프를 막지 않게 한다. 기동 계층은 전장을 모선 중심의 15채널 격자 지도로 래스터화하고,
FiLM 조건화를 갖춘 U-Net이 해면 전체에 대한 점수맵을 낸다. 각 방어정은 자신에게 허용된
유효 영역 안에서 점수가 높은 픽셀 두 개를 순차적으로 지목하고, 두 점을 잇는 선분에
그물벽을 놓는다. 연속 좌표를 회귀시키는 대신 점수맵 위의 픽셀을 지목하므로 행동공간이
격자로 규제되고, 적선 수가 달라져도 입력 형태가 유지된다. 학습은 가치망 없이 후보 행동을
직접 진행시켜 비교하는 그룹 상대 정책 최적화를 다중 에이전트로 확장하고, 한 척의 행동만
바꾸어 평가하는 반사실적 크레딧 배분으로 개별 기여를 분리한다. 전략 계층에서는 대규모
언어모델(LLM) 지휘관이 기하량이 미리 계산된 전장 상태를 받아 방어정별 담당 무리를 정하며,
조합 최적화에 의한 사후 보정을 두지 않아 배정의 책임이 전적으로 지휘관에게 있다.

\noindent
8개 조건 $\times$ 3개 공격 포메이션 $\times$ 30개 시드, 총 720 에피소드를 시드 대응표본으로
수집하여 두 계층의 기여를 분리 측정하였다. 기동 계층은 LLM 없이 단독으로 포획률을
$0.844 \to 0.886$ 으로 개선하였고(파상 포메이션에서 $+0.157$), 전략 계층은 지휘관 모델에
따라 결과가 갈렸다. 상용 클라우드 모델은 $+0.060$ 의 개선을 보인 반면 소형 로컬 모델
(Qwen2.5 7B)은 오히려 $-0.192$ 로 코드 휴리스틱보다 유의하게 나빴다. 10개 지표에 대한
$2\times2$ 상호작용 분해에서 9개가 신뢰구간에 0을 포함하여, 두 계층의 기여는 상승이 아니라
\emph{가산적}임을 확인하였다. 다중비교 보정에서는 포획률과 돌파 수가 유의성을 잃고 나머지
8개 지표는 유지되었다. 또한 포획률의 효과크기($\dz = 0.44$)가 측정한 10개 지표 중
최하위권이며, 총 선회량($-3.23$)과 포획당 그물 소모($-2.09$)에서 훨씬 큰 개선이 나타났다.
지휘관 계층의 성능은 배정 커버리지가 아니라 재계획 간 계획 안정성으로 설명되었다.
\end{abstract}
